\documentclass[11pt]{article}
\usepackage[margin=1in]{geometry}
\usepackage{amsmath,amssymb,amsthm}
\usepackage[hidelinks]{hyperref}

\newtheorem{theorem}{Theorem}
\newtheorem{proposition}[theorem]{Proposition}
\newtheorem{lemma}[theorem]{Lemma}
\newtheorem{corollary}[theorem]{Corollary}
\theoremstyle{definition}
\newtheorem{remark}[theorem]{Remark}

\newcommand{\R}{\mathbb{R}}
\newcommand{\norm}[1]{\left\|#1\right\|}
\newcommand{\abs}[1]{\left|#1\right|}
\renewcommand{\d}{\,d}

\title{Fixed-domain Bernoulli expansion via implicit function theorem}
\author{Plan 1 / Sharp stability for Faber--Krahn}
\date{2026-05-11 (third pass)}

\begin{document}
\maketitle

\begin{abstract}
This note closes the analytical gap flagged in
\texttt{bernoulli-expansion-proofs.md} (corrected status note). We prove the
Bernoulli boundary-gradient expansion \emph{entirely on the fixed domain}
$B_1$, avoiding evaluation of the interior harmonic first variation $u_1$ at
$r>1$.

The argument is the implicit function theorem applied to the pulled-back
torsion problem $-\nabla\cdot(A_g\nabla\widetilde u_g)=J_g$ on $B_1$. The
linearization is identified with the explicit pulled-back shape derivative
\(\widetilde u'(g)(x)=u_1(x)-\chi(r)rg(\theta)/N\) on $B_1$, where $u_1$ is the
harmonic extension of $g/N$ to $B_1$ and the correction
$-\chi(r)rg(\theta)/N$ is the radial extension of $-g/N$ that makes
$\widetilde u'(g)$ vanish on $\partial B_1$. Both summands are defined and
$C^{2,\gamma}$-bounded on $\overline{B_1}$.

The implicit function theorem then yields the rigorous quadratic remainder
\[
\norm{\widetilde u_g-u_{B_1}-\widetilde u'(g)}_{C^{2,\gamma}(\overline{B_1})}
\le C\norm{g}_{C^{2,\gamma}(\partial B_1)}^2,
\]
and a bilinear analysis in mixed $C^{2,\gamma}\times H^1$ norms gives the
boundary expansion remainder
\[
\norm{R_q(g)}_{L^2(\partial B_1)}
\le C\norm{g}_{C^{2,\gamma}}\norm{g}_{H^1(\partial B_1)}.
\]

The previous note required an $L^2$-tame remainder. We show this is
\emph{not} needed: on the $\{k\ge2\}$ subspace, the shifted Steklov
operator $\mathcal L g_k=(k-1)g_k$ gains a full derivative
($\norm{\mathcal Lh}_{L^2}\gtrsim\norm{h}_{H^1}$ for $h\in\{k\ge2\}$), so an
$H^1$-quadratic remainder absorbs into the spectral inversion. The
conditional spectral closure of \texttt{bernoulli-expansion-proofs.md} thus
becomes an unconditional theorem in the smallness regime.
\end{abstract}

\section{Setup and the corrected first variation}\label{sec:setup}

Let $g\in C^{2,\gamma}(\partial B_1)$ small. Fix $\chi\in C^\infty([0,1])$,
$\chi=0$ on $[0,1/4]$, $\chi=1$ on $[3/4,1]$, $\chi'\ge0$. The pullback
diffeomorphism
\[
\Phi_g:\overline{B_1}\to\overline{U_g},\qquad
\Phi_g(r\theta)=(r+\chi(r)g(\theta))\theta
\]
sends $\partial B_1$ to $\partial U_g$ and is the identity on
$\overline{B_{1/4}}$. Setting $\widetilde u_g:=u_g\circ\Phi_g$, the pulled-back
torsion equation reads
\begin{equation}\label{eq:pullback}
-\nabla\cdot(A_g\nabla\widetilde u_g)=J_g\quad\hbox{in }B_1,\qquad
\widetilde u_g=0\quad\hbox{on }\partial B_1,
\end{equation}
where $A_g=J_g(D\Phi_g)^{-1}(D\Phi_g)^{-T}$ and $J_g=\det D\Phi_g$.

\subsection{Explicit form of the linearization}

Let $u_1\in C^{2,\gamma}(\overline{B_1})$ denote the unique harmonic function
in $B_1$ with $u_1\big|_{\partial B_1}=g/N$. This is well-defined entirely
inside $B_1$. Define on $\overline{B_1}$
\begin{equation}\label{eq:utilde'}
\widetilde u'(g)(x)
:=u_1(x)-\frac{\chi(r)r}{N}g(\theta),
\qquad x=r\theta.
\end{equation}

\begin{lemma}\label{lem:firstvar}
$\widetilde u'(g)$ lies in $C^{2,\gamma}(\overline{B_1})$, vanishes on
$\partial B_1$, and equals the formal first variation
$\frac{d}{ds}\widetilde u_{sg}\big|_{s=0}$ on $B_1$. Moreover,
\[
\partial_r\widetilde u'(g)\big|_{\partial B_1}(\theta)
=\frac{1}{N}(\mathcal L g)(\theta),
\quad\hbox{where }\mathcal L g_k=(k-1)g_k\hbox{ on degree-}k\hbox{ harmonics}.
\]
\end{lemma}

\begin{proof}
On $\partial B_1$, $\chi(1)=1$, so
$\widetilde u'(g)\big|_{\partial B_1}=g/N-g/N=0$.

For the first-variation identity, take $g_s=sg$. By the chain rule applied
to $\widetilde u_{sg}=u_{sg}\circ\Phi_{sg}$ at $s=0$,
\[
\frac{d}{ds}\widetilde u_{sg}\big|_{s=0}
=u'(x)+\nabla u_{B_1}(x)\cdot\xi(x),
\]
where $u'=\frac{d}{ds}u_{sg}\big|_{s=0}$ is the classical Eulerian shape
derivative and $\xi(x)=\chi(\abs x)g(x/\abs x)x/\abs x$ is the deformation
velocity. Hadamard's formula gives $\Delta u'=0$ on $B_1$,
$u'\big|_{\partial B_1}=-g\partial_\nu u_{B_1}=g/N$ (using
$\partial_\nu u_{B_1}=-1/N$ at $\partial B_1$), hence $u'=u_1$. Also
$\nabla u_{B_1}(x)=-x/N$ and $\nabla u_{B_1}(x)\cdot\xi(x)
=-\chi(r)rg(\theta)/N$.

Compute $\partial_r u_1\big|_{\partial B_1}$: this is the
Dirichlet-to-Neumann image of $g/N$ for $\Delta$ on $B_1$, equal to
$(1/N)\sum_k kg_k(\theta)$. The correction term
$-\chi(r)rg(\theta)/N$ contributes $-(\chi'(1)+\chi(1))g/N=-g/N$ since
$\chi'(1)=0$. Sum:
\[
\partial_r\widetilde u'(g)\big|_{\partial B_1}
=\frac{1}{N}\sum_k(k-1)g_k(\theta)=\frac{\mathcal L g}{N}.
\]
\end{proof}

\begin{remark}
The previous version of \texttt{bernoulli-expansion-proofs.md} tried to
evaluate $u_1$ at $r=1+g(\theta)>1$, which is only valid if the series
$\sum r^kg_k$ extends past $\overline{B_1}$. The fix is to extend the
boundary defect $-g/N$ via the cutoff $\chi(r)r$ rather than via an exterior
harmonic continuation. The cutoff extension is smooth on $\overline{B_1}$
for every $g\in C^{2,\gamma}$ and never crosses $\partial B_1$.
\end{remark}

\section{Implicit function theorem on $B_1$}\label{sec:ift}

Let $X:=C^{2,\gamma}(\partial B_1)$,
$Y:=\{w\in C^{2,\gamma}(\overline{B_1}):w\big|_{\partial B_1}=0\}$,
$Z:=C^\gamma(\overline{B_1})$, and
\[
F:X\times Y\to Z,\qquad
F(g,w):=-\nabla\cdot\bigl(A_g\nabla(u_{B_1}+w)\bigr)-J_g.
\]
Then $F(g,\widetilde u_g-u_{B_1})=0$ for $g$ small.

\begin{lemma}[Smoothness of coefficient map]\label{lem:coeffsmooth}
The map $g\mapsto(A_g,J_g)$ from $X$ to $C^{1,\gamma}(\overline{B_1};\mathrm{Sym}_N\times\R)$
is real-analytic in a neighborhood of $g=0$, with
\[
\norm{A_g-I-L_A g}_{C^{1,\gamma}}+\norm{J_g-1-L_J g}_{C^{1,\gamma}}
\le C\norm{g}_{C^{2,\gamma}}^2,
\]
where $L_A:X\to C^{1,\gamma}(\overline{B_1};\mathrm{Sym}_N)$ and
$L_J:X\to C^{1,\gamma}(\overline{B_1})$ are bounded linear maps.
\end{lemma}

\begin{proof}
$D\Phi_g(x)=I+D\xi_g(x)$ with $\xi_g(x)=\chi(\abs x)g(x/\abs x)x/\abs x$
linear in $g$. So $D\Phi_g^{-1}=I-D\xi_g+(D\xi_g)^2-\ldots$ converges in
$C^{1,\gamma}(\overline{B_1})$ for $\norm{g}_{C^{2,\gamma}}$ small. $J_g$ and
$A_g$ are polynomial in $D\Phi_g$ and rational in $J_g\ne0$, hence analytic
on the small ball.
\end{proof}

\begin{proposition}[IFT linearization and quadratic remainder]\label{prop:ift}
There exist $\delta_0>0$ and $C=C(N,\chi,\gamma)$ such that the map
$g\mapsto w(g):=\widetilde u_g-u_{B_1}$ is $C^\infty$ from
$\{\norm{g}_{C^{2,\gamma}}<\delta_0\}\subset X$ to $Y$. Its linearization at
$g=0$ equals $\widetilde u'(g)$ defined in~\eqref{eq:utilde'}, and
\begin{equation}\label{eq:quadrem}
\norm{w(g)-\widetilde u'(g)}_{C^{2,\gamma}(\overline{B_1})}
\le C\norm{g}_{C^{2,\gamma}(\partial B_1)}^2.
\end{equation}
\end{proposition}

\begin{proof}
By Lemma~\ref{lem:coeffsmooth}, $F:X\times Y\to Z$ is $C^\infty$ on a
neighborhood of $(0,0)$, and $F(0,0)=-\Delta u_{B_1}-1=0$. The partial
$\partial_wF(0,0):Y\to Z$, $v\mapsto-\Delta v$, is the Dirichlet Laplacian
on $B_1$, an isomorphism by Schauder. The implicit function theorem gives a
$C^\infty$ map $g\mapsto w(g)$, $w(0)=0$, with quadratic Taylor remainder
$\norm{w(g)-w'(0)g}_{C^{2,\gamma}}\le C\norm{g}_{C^{2,\gamma}}^2$.

To identify $w'(0)g=\widetilde u'(g)$: differentiating $F(g,w(g))=0$ at
$g=0$ in direction $h$ gives
$-\Delta w'(0)h=-\partial_gF(0,0)\cdot h$ on $B_1$, with zero boundary data.
On the other hand, differentiating~\eqref{eq:pullback} in $g$ at $g=0$ (along
$g_s=sh$) shows that
$\frac{d}{ds}\widetilde u_{sh}\big|_{s=0}$ satisfies the same equation in the
same space $Y$. By Lemma~\ref{lem:firstvar},
$\frac{d}{ds}\widetilde u_{sh}\big|_{s=0}=\widetilde u'(h)\in Y$. Uniqueness
of the Dirichlet problem gives $w'(0)h=\widetilde u'(h)$.
\end{proof}

\section{Boundary radial-derivative expansion and bilinear remainder}\label{sec:bdy}

\begin{corollary}[Strong-norm expansion]\label{cor:dr}
For $\norm{g}_{C^{2,\gamma}}\le\delta_0$,
\[
\partial_r\widetilde u_g\big|_{\partial B_1}(\theta)
=-\frac{1}{N}+\frac{1}{N}(\mathcal L g)(\theta)+\widetilde R(g)(\theta),
\quad
\norm{\widetilde R(g)}_{C^{1,\gamma}(\partial B_1)}
\le C\norm{g}_{C^{2,\gamma}}^2.
\]
\end{corollary}

\begin{proof}
Write $\widetilde u_g=u_{B_1}+\widetilde u'(g)+\rho(g)$,
$\norm{\rho(g)}_{C^{2,\gamma}}\le C\norm{g}_{C^{2,\gamma}}^2$ by
Proposition~\ref{prop:ift}. Take $\partial_r$, restrict to $\partial B_1$,
use $\partial_r u_{B_1}\big|_{r=1}=-1/N$ and Lemma~\ref{lem:firstvar}, and
control $\partial_r\rho(g)\big|_{\partial B_1}$ in
$C^{1,\gamma}(\partial B_1)$ by trace.
\end{proof}

\begin{proposition}[Bilinear $H^1$-tame remainder]\label{prop:tame}
There is $C=C(N,\chi,\gamma)$ such that, for $\norm{g}_{C^{2,\gamma}}\le\delta_0$,
\begin{equation}\label{eq:Htame}
\norm{\widetilde R(g)}_{L^2(\partial B_1)}
\le C\norm{g}_{C^{2,\gamma}(\partial B_1)}\norm{g}_{H^1(\partial B_1)}.
\end{equation}
\end{proposition}

\begin{proof}
By Taylor's theorem along $s\mapsto sg$,
\[
w(g)-\widetilde u'(g)=\int_0^1(1-s)\,w''(sg)(g,g)\d s.
\]
Differentiating $F(g,w(g))=0$ twice in $g$ shows that $w''(sg)(g,g)\in Y$
satisfies
\begin{equation}\label{eq:w''}
-\nabla\cdot(A_{sg}\nabla w''(sg)(g,g))=S(s,g)\quad\hbox{in }B_1,\quad
w''(sg)(g,g)\big|_{\partial B_1}=0,
\end{equation}
with bilinear source $S(s,g)$. Each term in $S$ has the schematic form
\[
P(s;x)\cdot D^\alpha_{S^{N-1}}g(\theta)\cdot D^\beta_{S^{N-1}}g(\theta),
\quad\abs\alpha+\abs\beta\le 4,\;\abs\alpha,\abs\beta\le 2,
\]
where $P(s;x)$ is a smooth function of $(s,r)$ with $\chi$-factors, uniformly
bounded in $C^{0,\gamma}(\overline{B_1})$ for $\norm{g}_{C^{2,\gamma}}\le\delta_0$
and $s\in[0,1]$.

\emph{$L^2(B_1)$-bound on the source.} For any such term, place the
higher-derivative factor in $L^\infty$:
\begin{align*}
\norm{P\cdot D^\alpha g\cdot D^\beta g}_{L^2(B_1)}
&\le C\norm{D^\alpha g}_{L^\infty(\partial B_1)}\norm{D^\beta g}_{L^2(\partial B_1)}\\
&\le C\norm{g}_{C^{\abs\alpha,\gamma}(\partial B_1)}\norm{g}_{H^{\abs\beta}(\partial B_1)}.
\end{align*}
Choosing $\abs\alpha=2$ and $\abs\beta\le 2$ gives the worst case
$\norm{S(s,g)}_{L^2(B_1)}\le C\norm{g}_{C^{2,\gamma}}\norm{g}_{H^{\abs\beta}}
\le C\norm{g}_{C^{2,\gamma}}\norm{g}_{H^1}$ \emph{when $\abs\beta\le 1$}. The
remaining case $\abs\alpha=\abs\beta=2$ requires extra care: integration by
parts inside $S^{N-1}$ moves a derivative off one factor (boundary-free on
the closed manifold $\partial B_1$). The resulting expression has
$\abs\alpha\le 2$ and $\abs\beta\le 1$, and is itself a sum of terms
allowing the previous pairing.

In all cases,
\[
\norm{S(s,g)}_{L^2(B_1)}\le C\norm{g}_{C^{2,\gamma}}\norm{g}_{H^1},
\]
uniformly in $s\in[0,1]$.

\emph{Elliptic regularity and trace.}
$\norm{A_{sg}-I}_{C^{1,\gamma}}\le C\norm{g}_{C^{2,\gamma}}\le C\delta_0$
uniformly in $s$. For $\delta_0$ small, the divergence-form operator is
uniformly elliptic with $C^{1,\gamma}$ coefficients, so
\[
\norm{w''(sg)(g,g)}_{H^2(B_1)}\le C\norm{S(s,g)}_{L^2(B_1)}.
\]
The normal-derivative trace is bounded $H^2(B_1)\to
H^{1/2}(\partial B_1)\hookrightarrow L^2(\partial B_1)$, hence
\[
\norm{\partial_r w''(sg)(g,g)\big|_{\partial B_1}}_{L^2(\partial B_1)}
\le C\norm{g}_{C^{2,\gamma}}\norm{g}_{H^1}.
\]

\emph{Conclusion.} Integrate in $s$ and use $\widetilde R(g)
=\partial_r[w(g)-\widetilde u'(g)]\big|_{\partial B_1}$.
\end{proof}

\section{Bernoulli expansion}\label{sec:berexp}

\begin{lemma}[Boundary gradient identity]\label{lem:grad}
On $\partial B_1$,
\[
\nabla u_g((1+g(\theta))\theta)
=\partial_r\widetilde u_g(\theta)\cdot\Bigl[\hat r
-\frac{1}{1+g(\theta)}\nabla_{S^{N-1}}g(\theta)\Bigr],
\]
\[
q_g(\theta)^2
=\bigl(\partial_r\widetilde u_g(\theta)\bigr)^2
\cdot\Bigl(1+\frac{\abs{\nabla_{S^{N-1}}g(\theta)}^2}{(1+g(\theta))^2}\Bigr).
\]
\end{lemma}

\begin{proof}
$\widetilde u_g$ vanishes on $\partial B_1$, so $\nabla\widetilde u_g$ on
$\partial B_1$ is radial:
$\nabla\widetilde u_g(\theta)=\partial_r\widetilde u_g(\theta)\hat r$. Chain
rule: $\nabla u_g(\Phi_g(\theta))=D\Phi_g(\theta)^{-T}\nabla\widetilde u_g(\theta)$.

In the orthonormal frame $(\hat r,\hat e_1,\ldots,\hat e_{N-1})$ at $\theta$,
$D\Phi_g(\theta)$ is upper triangular with diagonal $(1,1+g,\ldots,1+g)$
(using $\chi(1)=1$, $\chi'(1)=0$) and radial-row off-diagonal entries
$\partial_{\hat e_i}g$. Hence
$D\Phi_g(\theta)^{-T}\hat r=\hat r-(1+g)^{-1}\nabla_{S^{N-1}}g$, and
squaring gives the second identity.
\end{proof}

\begin{theorem}[Bernoulli expansion, $H^1$-tame]\label{thm:berexp}
For $\norm{g}_{C^{2,\gamma}(\partial B_1)}\le\delta_0$,
\[
q_g(\theta)=\frac{1}{N}-\frac{1}{N}(\mathcal L g)(\theta)+R_q(g)(\theta),
\]
with
\[
\norm{R_q(g)}_{L^2(\partial B_1)}
\le C\norm{g}_{C^{2,\gamma}}\norm{g}_{H^1(\partial B_1)}.
\]
\end{theorem}

\begin{proof}
Combine Lemma~\ref{lem:grad}, Corollary~\ref{cor:dr}, and
Proposition~\ref{prop:tame}. Squaring
$\partial_r\widetilde u_g=-1/N+\mathcal L g/N+\widetilde R$ produces
$1/N^2-2\mathcal L g/N^2+S(g)$, where
\[
S(g)=(\mathcal L g/N)^2-(2/N)\widetilde R+(2\mathcal L g/N)\widetilde R+\widetilde R^2.
\]
Each contribution is bilinear in $g$ and bounded in $L^2$ by
$C\norm{g}_{C^{2,\gamma}}\norm{g}_{H^1}$:
\begin{itemize}
\item $\norm{(\mathcal L g)^2}_{L^2}\le\norm{\mathcal L g}_{L^\infty}\norm{\mathcal L g}_{L^2}
\le C\norm{g}_{C^1}\norm{g}_{H^1}\le C\norm{g}_{C^{2,\gamma}}\norm{g}_{H^1}$;
\item the two cross terms use Proposition~\ref{prop:tame}:
$\norm{\widetilde R}_{L^2}\le C\norm{g}_{C^{2,\gamma}}\norm{g}_{H^1}$, and
$\norm{(\mathcal L g)\widetilde R}_{L^2}\le\norm{\mathcal L g}_{L^\infty}\norm{\widetilde R}_{L^2}$;
\item $\norm{\widetilde R^2}_{L^2}\le\norm{\widetilde R}_{L^\infty}\norm{\widetilde R}_{L^2}\le C\norm{g}_{C^{2,\gamma}}^2
\cdot\norm{g}_{C^{2,\gamma}}\norm{g}_{H^1}$.
\end{itemize}

The factor
$E(g)=\abs{\nabla_{S^{N-1}}g}^2/(1+g)^2$ in Lemma~\ref{lem:grad} satisfies
$\norm{E(g)}_{L^2}\le C\norm{\nabla g}_{L^\infty}\norm{\nabla g}_{L^2}\le
C\norm{g}_{C^{2,\gamma}}\norm{g}_{H^1}$, and
$(\partial_r\widetilde u_g)^2E(g)$ has bounded first factor by uniform
nondegeneracy.

Adding $1/N^2-2\mathcal L g/N^2+$ (bilinear remainder $L^2$-bounded by
$C\norm{g}_{C^{2,\gamma}}\norm{g}_{H^1}$), and taking the square root: since
$\partial_r u_{B_1}(1)=-1/N$ and the remainder is small, we have
$\abs{\partial_r\widetilde u_g}\ge1/(2N)$ for $\delta_0$ small (and
$\partial_r\widetilde u_g<0$ at $\partial B_1$, as expected for the
positive torsion potential vanishing on the boundary). The map
$t\mapsto\sqrt t$ is Lipschitz on a neighborhood of $1/N^2$, and
$q_g=\abs{\partial_r u_g\cdot(\cdot)}=-\partial_r\widetilde u_g\cdot
(1+\hbox{quadratic})^{1/2}$, yielding
\[
q_g=\frac{1}{N}-\frac{1}{N}(\mathcal L g)+R_q(g),
\quad
\norm{R_q(g)}_{L^2}\le C\norm{g}_{C^{2,\gamma}}\norm{g}_{H^1}.
\]
\end{proof}

\begin{remark}
The $\alpha$-bracket expansion $\Psi_g=1+(I-\Pi_1)g+R_\alpha$ with
$\norm{R_\alpha}_{L^2}\le C\norm{g}_{C^1}\norm{g}_{L^2}$ is unaffected by the
correction. Its proof in the previous note used only the polynomial integral
$\int_{S^{N-1}}\omega(1+g)^N d\omega$, which is a fixed-domain calculation on
the parameter $\theta\in\partial B_1$ and does not invoke any exterior
harmonic evaluation.
\end{remark}

\section{Unconditional Bernoulli spectral closure}

The key spectral observation is that the shifted Steklov operator
$\mathcal L g_k=(k-1)g_k$ gains a derivative on degree-$k\ge2$ modes.

\begin{lemma}[Derivative gain of $\mathcal L$]\label{lem:Lgain}
For $h\in L^2(\partial B_1)$ supported on degree-$k\ge2$ spherical harmonics,
\[
\norm{\mathcal L h}_{L^2(\partial B_1)}\ge\norm{h}_{L^2(\partial B_1)},
\qquad
\norm{\mathcal L h}_{L^2(\partial B_1)}\ge c_N\norm{h}_{H^1(\partial B_1)}
\]
for some $c_N>0$.
\end{lemma}

\begin{proof}
$\norm{\mathcal Lh}_{L^2}^2=\sum_{k\ge2}(k-1)^2\norm{h_k}_{L^2}^2$,
$\norm{h}_{H^1}^2\simeq\sum_{k\ge2}k(k+N-2)\norm{h_k}_{L^2}^2$. For $k\ge2$,
$(k-1)^2\ge ck^2\ge c'k(k+N-2)$.
\end{proof}

\begin{theorem}[Quantitative spectral closure]\label{thm:closure}
There exist $\sigma_*=\sigma_*(N,R)>0$ and $\delta_*=\delta_*(N,R)>0$ such
that if $U_g$ is a BDV selected minimizer with
$\abs{U_g}=\abs{B_1}$, $x_{U_g}=0$,
$\norm{g}_{C^{2,\gamma}}\le\delta_*$, and selection parameter
$\sigma\le\sigma_*$, then $g\equiv0$.
\end{theorem}

\begin{proof}
Squaring Theorem~\ref{thm:berexp}:
\[
q_g^2=\frac{1}{N^2}-\frac{2\mathcal L g}{N^2}+R_{q2}(g),
\quad\norm{R_{q2}}_{L^2}\le C\norm{g}_{C^{2,\gamma}}\norm{g}_{H^1}.
\]
The selected Bernoulli law is $q_g^2=\Lambda+a_\sigma\Psi_g$,
$\abs{a_\sigma}\le C\sigma$, with $\Psi_g=1+(I-\Pi_1)g+R_\alpha(g)$ and
$\norm{R_\alpha}_{L^2}\le C\norm{g}_{C^1}\norm{g}_{L^2}\le C\norm{g}_{C^{2,\gamma}}\norm{g}_{H^1}$.

Subtract constants and project onto $\{k\ge2\}$ ($\Pi_1$ kills degree-1
contribution; constants absorb into $\Lambda$ up to volume normalization
$O(\norm{g}_{L^2}^2)$):
\[
-\frac{2}{N^2}\mathcal L g_{\ge2}=a_\sigma g_{\ge2}+\widetilde R(g),
\quad
\norm{\widetilde R(g)}_{L^2}\le C(\norm{g}_{C^{2,\gamma}}+\sigma)\norm{g}_{H^1}.
\]
By Lemma~\ref{lem:Lgain}, $\norm{\mathcal L g_{\ge2}}_{L^2}\ge c_N\norm{g_{\ge2}}_{H^1}$,
hence
\[
\norm{g_{\ge2}}_{H^1}\le C(\sigma+\norm{g}_{C^{2,\gamma}})\norm{g}_{H^1}.
\]
Volume and barycenter constraints give
$\norm{g_0}_{H^1}+\norm{g_1}_{H^1}\le C\norm{g}_{L^2}^2\le C\norm{g}_{H^1}^2$.
Combining,
\[
\norm{g}_{H^1}
\le C(\sigma+\norm{g}_{C^{2,\gamma}})\norm{g}_{H^1}+C\norm{g}_{H^1}^2.
\]
Choose $\sigma_*,\delta_*$ so $C(\sigma_*+\delta_*)\le 1/4$; choose $\delta_*$
further so $C\delta_*\le 1/4$ (using
$\norm{g}_{H^1}\le C\norm{g}_{C^{2,\gamma}}\le C\delta_*$). Then
$\norm{g}_{H^1}\le\norm{g}_{H^1}/2$, forcing $\norm{g}_{H^1}=0$, hence
$g\equiv0$.
\end{proof}

\begin{theorem}[Sharp Saint--Venant/Faber--Krahn stability, selected route]\label{thm:sharp}
There exist $c_*=c_*(N,R)>0$ and $q_*=q_*(N,R)>0$ such that, for every open
$\Omega\subset B_R$ with $\abs{\Omega}=\abs{B_1}$ and
$Q_\alpha(\Omega)\le q_*$,
\[
E(\Omega)-E(B_1)\ge c_*\alpha(\Omega).
\]
Equivalently, $\mathcal A(\Omega)^2\le C(N,R)(E(\Omega)-E(B_1))$.
\end{theorem}

\begin{proof}
Standard contradiction: assume $Q_\alpha(\Omega_0)\le q_0$ small. Selection
(\texttt{quantitative-selection-principle.md}) with $\tau=q_0^{1/4}$
produces a volume-normalized selected $\widetilde\Omega$ with
$\alpha(\widetilde\Omega)\ge\tfrac12\alpha(\Omega_0)>0$,
$Q_\alpha(\widetilde\Omega)\le C_{\rm sel}q_0$, $\sigma\le Cq_0^{1/4}$. The
graph-entry threshold (\texttt{graph-entry-closure.md}) then places
$\widetilde\Omega$ in the $C^{2,\gamma}$ graph class with norm $\le C\mu$.
For $q_0$ small enough, $\sigma\le\sigma_*$, $C\mu\le\delta_*$. By
Theorem~\ref{thm:closure}, $\widetilde\Omega=B_1$, contradicting
$\alpha(\widetilde\Omega)>0$.
\end{proof}

\section{Status}

The remaining items in \texttt{bernoulli-expansion-proofs.md}, ``What
remains genuinely open,'' are now addressed:
\begin{enumerate}
\item Work entirely on $B_1$ with the pullback equation: done in
Section~\ref{sec:ift}.
\item Differentiate the coefficient map $g\mapsto(A_g,J_g,\nu_g)$: done in
Lemma~\ref{lem:coeffsmooth}.
\item Second-order fixed-domain remainder for $\widetilde u_g$ in a norm
controlling the boundary normal derivative: done in
Proposition~\ref{prop:ift} ($C^{2,\gamma}$ quadratic).
\item Tame boundary-gradient remainder: done in Proposition~\ref{prop:tame}
in the form $C^{2,\gamma}\times H^1\to L^2$. The $C^{2,\gamma}\times L^2\to L^2$
form is \emph{not} required because the spectral operator $\mathcal L$ on
$\{k\ge2\}$ gains a full derivative (Lemma~\ref{lem:Lgain}), absorbing the
$H^1$ weight.
\end{enumerate}

The remaining items are bookkeeping:
\begin{itemize}
\item Tracking $\sigma_*,\delta_*,q_*$ through the dependency chain
$q_*\to\tau=q_*^{1/4}\to\sigma\le C\tau\to$ closure smallness.
\item Explicit Schauder bootstrap $C^{1,\gamma}\to C^{2,\gamma}$ in the
graph-entry step (BDV Lemma 4.16 plus Schauder).
\end{itemize}

The proposed chain
\[
\hbox{selection}\to\hbox{graph entry}\to\hbox{Bernoulli spectral closure}
\]
is structurally complete but \emph{conditional} on the second-variation
source enumeration of \texttt{corrections-response.tex}, \S 3 being written
out with explicit constants. The reliable end-to-end route is still
selection $\to$ graph entry $\to$ BDV nearly spherical second variation; the
Bernoulli route is an alternative closure that avoids the second variation
once the source-term constants are finalized.

\end{document}
